% !BIB program = none
%% ============================================================
%%  The Fluid Mechanics and Policy Case for Bidet Adoption
%%  in the United States: A Quantitative Analysis
%%
%%  REVISION (30 May 2026):
%%    - Standoff narrative reconciled with companion optimization
%%      paper (bidet_optimization_1.tex): the previous
%%      ``momentum-conservation insensitivity to standoff'' claim
%%      is retracted and replaced with the developed-jet decay
%%      treatment of the companion paper.
%%    - Notation clash a (nozzle radius) vs a (stagnation strain
%%      rate) resolved by renaming the strain rate to \dot epsilon_s.
%%    - r_clean evaluation corrected: previous 16 mm was a numerical
%%      slip; the correct value for the stated parameters is
%%      approximately 26 mm.
%%    - Reynolds-number convention unified: Re_d everywhere.
%%    - Angular-coverage equation derivation cleaned up.
%%    - Hydraulic-jump treatment replaced with Bohr-Dimon-Putkaradze
%%      shallow-water-shock scaling (Watson's original gravity
%%      formula is dimensionally inconsistent with the stated
%%      parameter values).
%%    - Hussein, Capp & George 1994 added for B constant;
%%      Bohr et al. 1993 / Bush & Aristoff 2003 added for jump.
%%    - Companion-paper cross-references added; verification roadmap
%%      cross-referenced.
%% ============================================================
\documentclass[10pt,letterpaper]{article}
\usepackage[margin=1in]{geometry}
\usepackage{natbib}

% ---- Packages -----------------------------------------------
\usepackage{graphicx}
\usepackage{booktabs}
\usepackage{array}
\usepackage{multirow}
\usepackage{xcolor}

, group-minimum-digits=4}
\usepackage{enumitem}
\usepackage{tikz}
\usetikzlibrary{arrows.meta, decorations.pathmorphing, patterns}
\usepackage[colorlinks=true, linkcolor=blue!60!black,
            citecolor=blue!60!black, urlcolor=blue!60!black]{hyperref}

% ---- Custom commands ----------------------------------------
\newcommand{\Rey}{\mathrm{Re}}
\newcommand{\We}{\mathrm{We}}
\newcommand{\Pe}{\mathrm{Pe}}
\newcommand{\Bi}{\mathrm{Bi}}
\newcommand{\Go}{\mathrm{G\ddot{o}}}
\newcommand{\order}[1]{\mathcal{O}\!\left(#1\right)}
\newcommand{\half}{\tfrac{1}{2}}

% =============================================================
\begin{document}

\title{On the Superiority of Hydrodynamic Perianal Cleansing: A
       Quantitative Framework in Fluid Mechanics, Tribology, and
       Environmental Economics for Bidet Adoption Policy in the
       United States}

\author{Charlie Burton}
%\email{cpb@u.northwestern.edu}
%\affiliation{Department of Physics and Astronomy, Northwestern University}

\date{\today}

\begin{abstract}
The United States remains an international outlier in its near-universal
reliance on dry toilet paper (TP) for perianal hygiene, consuming
141~rolls/ person/ year, the highest per-capita rate on
Earth, while countries such as Japan have achieved 80~\% household
bidet penetration. This paper constructs a quantitative argument
for bidet adoption in American residential and commercial buildings across
four domains. Using Watson's (1964) radial wall-jet solution augmented by a
curvature-perturbation expansion, we model the nozzle as a circular
free jet at exit Reynolds number $\Rey_d \approx 10{,}000$--$18{,}000$
that transitions to a wall jet of laminar boundary-layer character at
the cleaning radius, and derive wall shear stress
distributions $\tau_w(r)$; cleaning-zone values are estimated at
$\tau_w \approx 50$--$500~Pa$ in the laminar lower-bound limit
of Watson Region~II (Phares et al.~2000 wall-jet correlations imply
that the true turbulent value is several-fold higher), and either
estimate exceeds the critical detachment threshold for fecal residue
($\tau_\mathrm{crit} \approx 1$--$10~Pa$) by one to two orders
of magnitude. A P\'{e}clet-number analysis ($\Pe \sim 10^9$)
confirms that advective transport in the spreading water film eliminates
dislodged material in a single pass, whereas dry wiping redistributes
rather than removes debris. Tribological analysis of toilet paper on moist
perianal skin yields contact shear stresses of 1--30~kPa, two to
three orders of magnitude above the bidet's hydrodynamic shear, providing
a mechanistic basis for stratum corneum disruption and the associated
clinical spectrum of pruritus ani, anal fissures, and hemorrhoidal
exacerbation. An environmental lifecycle and economic analysis shows that
national bidet adoption would eliminate the harvest of 15--31 million trees
annually, reduce embedded water consumption by $\sim98~\%$, and pay
back household installation costs within 3--15 months. A targeted policy
proposal modeled on the EPA WaterSense rebate framework is presented.
\end{abstract}

\maketitle

% =============================================================
\section{Introduction}
\label{sec:intro}
% =============================================================

The act of perianal cleansing following defecation is performed by nearly
every human being on the planet multiple times daily, yet the dominant
American method (dry mechanical wiping with cellulose paper) has received
almost no formal engineering scrutiny. The bidet employs a directed water
jet to accomplish the same task through different physics: convective mass
transfer rather than frictional contact, hydrodynamic shear rather than
tribological abrasion, and advective removal rather than smearing. These
differences are quantitative, spanning two to three orders of magnitude in
the mechanical stress delivered to perianal tissue.

The public health stakes are substantial. Hemorrhoidal disease affects an
estimated 10 million Americans annually, generating \$770 million in
documented employer-insured healthcare costs \citep{Peery2014}. Pruritus
ani has a population prevalence of 1--5\%, with mechanical irritation from
wiping identified as a primary or exacerbating factor in a significant
fraction of cases \citep{Dasan1999}. The environmental costs are no less
significant: American toilet paper consumption accounts for approximately
20\% of global supply \citep{ScientificAmerican2009}, with 80\% manufactured
from virgin forest fiber \citep{NRDC2020}.


% =============================================================
\section{Anatomical Geometry: The Hemi-Cylinder Model}
\label{sec:geometry}
% =============================================================

\subsection{Geometric Idealization}

The perianal anatomy presents a concave ellipsoidal cleft whose local geometry
near the anal verge is well approximated, for the purpose of fluid mechanical
analysis, by a hemi-cylindrical surface. We define the \emph{effective cleft
radius} $R$ as the radius of curvature of the inter-gluteal fold at the level
of the anus, measured in the sagittal plane. From anthropometric data,
$R \approx 15$--$25~mm$ for the seated adult posture relevant to bidet
use.

We introduce cylindrical coordinates $(r, \phi, z)$ with $z$ directed along
the jet axis (vertically downward from the nozzle), the apex of the concave
surface located at the origin, and $r$ measuring radial distance from the
stagnation point on the surface. The concave surface profile is described by
\begin{equation}
    z_s(r) = \frac{r^2}{2R} + \mathcal{O}\!\left(\frac{r^4}{R^3}\right),
    \qquad r \ll R,
    \label{eq:surfaceprofile}
\end{equation}
i.e., a paraboloid of revolution to leading order. The flat-plate limit
corresponds to $R \to \infty$, and curvature effects enter at
$\mathcal{O}(r^2/R^2)$.

The \emph{curvature parameter} is defined as
\begin{equation}
    \epsilon \equiv \frac{d}{R},
    \label{eq:curvparam}
\end{equation}
where $d$ is the nozzle diameter. For typical values $d \approx 1$--$2~mm$
and $R \approx 20\,\mathrm{mm}$, we have $\epsilon \approx 0.05$--$0.10$,
confirming that curvature is a small but non-negligible perturbation. The
analysis of Sections~\ref{sec:fluid}--\ref{sec:mass} proceeds by (i)
solving the flat-plate impinging jet problem exactly, and (ii) incorporating
curvature effects via a systematic perturbation expansion in $\epsilon$.

\subsection{Debris Distribution and the Cleaning Criterion}

Fecal residue is deposited on the concave surface as a thin layer of
shear-thinning viscoplastic material. Let $\Sigma(r)$ denote the areal
density of adhered material (kg/m$^2$) as a function of radial position.
The cleaning criterion requires that the wall shear stress delivered by the
bidet jet exceed the adhesive yield stress of the fecal layer at every point:
\begin{equation}
    \tau_w(r) \geq \tau_\mathrm{crit}(r) \quad \forall\, r \in [0, r_\mathrm{max}],
    \label{eq:cleaningcriterion}
\end{equation}
where $r_\mathrm{max}$ is the maximum radial extent of contamination (estimated
at 10--20~mm based on anatomical considerations).

The effective angular cleaning coverage on the hemi-cylinder is bounded
by the polar angle $\theta_\mathrm{max}$ at which the impinging
normal-velocity component, taken as $V_n(\theta) = V_0 \cos\theta$ for
oblique impingement on a locally tangent plane
\citep{BeitelmalEtAl2000}, falls below the value that, when inserted
into Watson's Region~II shear law~\eqref{eq:tauII}, satisfies
$\tau_w(r_\mathrm{max}) = \tau_\mathrm{crit}$. Because $\tau_w$ scales as
$V_n^3$ at fixed nozzle geometry,
\begin{equation}
    \cos^3\theta_\mathrm{max}
    \;=\; \frac{\tau_\mathrm{crit}\,\nu\,r_\mathrm{max}^5}
              {c_2\,\rho\,V_0^3\,a^6}\;\equiv\;\chi,
    \label{eq:angularcoverage}
\end{equation}
so that $\theta_\mathrm{max} = \arccos(\chi^{1/3})$. For the nominal
$\tau_\mathrm{crit} \approx 5~Pa$, $V_0 \approx 11.2~m/s$,
$a = 0.65~mm$, $r_\mathrm{max} = 15~mm$, $\nu = e-6~m^2/s$,
equation~\eqref{eq:angularcoverage} gives $\chi \approx 0.005$ and
$\theta_\mathrm{max} \approx 80^\circ$, indicating that effective
cleaning extends over essentially the entire hemi-cylindrical
surface. The result is, however, sensitive to the Region~II
laminar conservatism (Sec.~\ref{subsec:watson}): a turbulent
correction enlarges the cleaning angle further.

% =============================================================
\section{Fluid Mechanics of the Impinging Bidet Jet}
\label{sec:fluid}
% =============================================================

\subsection{Nozzle Conditions and Free-Jet Parameters}
\label{subsec:nozzle}

A modern bidet nozzle delivers a coherent circular jet through a circular
orifice of diameter $d = 1$--$2~mm$ at a supply gauge pressure
$\Delta P = 0.05$--$0.1~MPa$. The jet exit velocity follows from the
Torricelli relation with a velocity coefficient $C_v \approx 0.97$:
\begin{equation}
    V_0 = C_v \sqrt{\frac{2\,\Delta P}{\rho}},
    \label{eq:torricelli}
\end{equation}
yielding $V_0 \approx 8$--$14~m/s$ across the operating envelope.
The volumetric flow rate and mass-flux Reynolds number are, respectively,
\begin{equation}
    Q = \frac{\pi d^2}{4} V_0, \qquad
    \Rey_d = \frac{V_0\,d}{\nu},
    \label{eq:reynolds}
\end{equation}
with $\nu = 8.9 \times 10^{-7}\,\mathrm{m^2/s}$ for water at
$T = 25~^\circ\mathrm{C}$. For $V_0 = 10~m/s$ and $d = 1.3~mm$:
\begin{equation}
    \Rey_d = \frac{(10)(1.3 \times 10^{-3})}{8.9 \times 10^{-7}}
    \approx 14{,}600,
\end{equation}
placing the exit flow in the fully turbulent regime ($\Rey_d > 3{,}000$
for a pipe, $>4{,}000$ for a free jet \citep{ZuckermanLior2006}).

\subsubsection{Jet Coherence: Weber Number Analysis}

The Weber number governs whether surface tension can break the jet into
droplets via Rayleigh--Plateau instability:
\begin{equation}
    \We = \frac{\rho V_0^2\,d}{\sigma},
    \label{eq:weber}
\end{equation}
where $\sigma = 0.072~N/m$ for water--air at $25~^\circ\mathrm{C}$.
For the nominal parameters, $\We \approx 1{,}400$--$2{,}800$, indicating
that inertia overwhelmingly dominates surface tension. The axisymmetric
Rayleigh breakup length, accounting for the turbulent perturbation spectrum
\citep{Lefebvre1988}, scales as
\begin{equation}
    L_\mathrm{break} \approx 12\,\We^{1/2}\,d \approx 450\,\mathrm{mm}.
    \label{eq:breakup}
\end{equation}
Since the nozzle-to-skin standoff $H$ is typically $H \approx 50$--$150~mm$,
the condition $L_\mathrm{break} \gg H$ is satisfied comfortably. The jet
therefore arrives at the skin surface as a continuous liquid column of
nearly uniform velocity $V_0$, not as a dispersed spray. A spray carries
less momentum per unit area and generates lower wall shear stress.

\subsubsection{Free-Jet Velocity Decay}

Beyond the potential core, which extends to $z/d \approx 4$--$6$
\citep{HusseinCappGeorge1994,PopeBook}, the centreline velocity of a
self-similar axisymmetric turbulent jet decays as
\begin{equation}
    \frac{V_c(z)}{V_0} \;=\; \frac{B}{z/d},
    \qquad B \approx 5.8\text{--}6.1,
    \label{eq:jetdecay}
\end{equation}
\citep{HusseinCappGeorge1994,PopeBook}, with a self-similar Gaussian
radial profile of half-width $r_{1/2}(z) \approx 0.094\,z$. We take
$B = 6.0$ in numerical estimates that follow.

For a standoff of $H = 80~mm$ and $d = 1.3~mm$ we have
$H/d \approx 62$, so the impinging centreline velocity is
$V_c(H) = V_0\,B/(H/d) \approx 1.1~m/s$, a tenfold drop from $V_0$.
A naive momentum-conservation argument that treats the wall jet as
driven by the conserved axial momentum flux $J = \rho \pi d^2 V_0^2/4$
appears to make the impingement insensitive to this decay, but \emph{the
input to Watson's similarity solution is the local impingement-velocity
distribution, not the total flux}. Once the free jet has spread to a
Gaussian footprint of width $\sim 0.094 H$ at impingement, substituting
the spread profile and the decayed centreline into Watson's
Region~II shear (Sec.~\ref{subsec:watson}) yields a wall shear that
falls as $\tau_w \propto (d/H)^{2}$, equivalent to a minimum-flow penalty
$Q^* \propto (H/d)^{2/3}$ relative to $H/d \to 1$; the companion
optimization paper~\citep{BurtonOptPaper} works the algebra in detail.
For the conventional seated-user geometry ($H/d \approx 60$) versus a
close-deploy nozzle ($H/d \approx 6$) the penalty is
$\sim 10^{2/3} \approx 4.6\times$, identifying standoff reduction as the
dominant water-conservation lever rather than a minor correction.
\emph{We retract the ``relative insensitivity to jet decay'' claim of
the previous revision.}

\subsection{Impingement Zone: Stagnation Flow Analysis}
\label{subsec:stagnation}

In the immediate vicinity of the stagnation point ($r \lesssim d/2$), the
flow is described by axisymmetric stagnation-point theory. Along the
stagnation streamline, the continuity equation requires the radial velocity
to grow linearly with $r$:
\begin{equation}
    u_r(r, 0) \;=\; \dot\varepsilon_s\,r + \mathcal{O}(r^3),
    \label{eq:radvel}
\end{equation}
where $\dot\varepsilon_s$ is the \emph{strain rate} at the stagnation
point (renamed from $a$ in earlier drafts to avoid collision with the
nozzle radius $a \equiv d/2$). For an impinging jet, Schlichting's
boundary-layer analysis gives $\dot\varepsilon_s \approx V_c(H)/
l_\mathrm{imp}$, with the impingement length scale $l_\mathrm{imp}
\sim d/2$. The Bernoulli equation applied along the stagnation
streamline yields the \emph{stagnation pressure}:
\begin{equation}
    P_s = P_\infty + \frac{1}{2}\rho V_c^2(H).
    \label{eq:stagpress}
\end{equation}

For the typical seated-user standoff $H/d \approx 60$, the centreline
velocity at impingement decays as $V_c(H) = B V_0 d/H$ with $B \approx 6.3$
\citep{PopeBook}, giving $V_c \approx 1.2~m/s$ and a stagnation
overpressure of only
\begin{equation}
    P_s - P_\infty \approx \tfrac{1}{2}(1000)(1.2)^2 \approx 0.7~kPa
    \quad \text{(typical $H/d \sim 60$)}.
    \label{eq:stagval}
\end{equation}
For a close-deploy nozzle ($H/d \lesssim 6$, $V_c \approx V_0$), the
short-standoff value is
\begin{equation}
    P_s - P_\infty \approx \tfrac{1}{2}(1000)(10)^2 = 50~kPa
    \quad \text{(close-deploy $H/d \lesssim 6$)}.
    \label{eq:stagshort}
\end{equation}
In either regime, the stagnation pressure is compressive (normal to the
surface), acts over a very small contact area
$A_\mathrm{stag} \sim \pi d^2/4 \approx 1.3~mm^2$, and produces no
frictional shear. The stagnation region therefore poses no abrasive risk to
the skin; the loading is purely hydrostatic.

The pressure distribution along the surface decays from the stagnation value
as (from Bernoulli and the radial acceleration in equation~\ref{eq:radvel})
\begin{equation}
    P(r) \;=\; P_s - \tfrac{1}{2}\rho\,\dot\varepsilon_s^2\,r^2.
    \label{eq:pressuredist}
\end{equation}
The radial pressure gradient
$\partial P/\partial r = -\rho\,\dot\varepsilon_s^2\,r < 0$
drives the flow outward into the spreading wall jet.

\subsection{Wall Jet Region: Watson's Solution for the Flat-Plate Baseline}
\label{subsec:watson}

Beyond the stagnation zone ($r > d/2$), the flow transitions into a thin
radially spreading \emph{wall jet} governed by the Prandtl boundary-layer
equations. Following Watson (1964) \citep{Watson1964}, we consider the
axisymmetric spreading of a circular liquid jet on a horizontal plate under
the assumptions: (i) the outer flow is inviscid with uniform surface velocity
$U_s \approx V_0$; (ii) a viscous boundary layer of thickness $\delta(r)$
grows from the plate upward; and (iii) the pressure is atmospheric throughout
(no streamwise adverse pressure gradient in the outer film, as verified
by the decay of equation~\ref{eq:pressuredist}).

The governing equations in the boundary layer are
\begin{equation}
    \frac{\partial(r\,u_r)}{\partial r} + \frac{\partial(r\,u_z)}{\partial z}
    = 0,
    \qquad
    u_r\,\frac{\partial u_r}{\partial r} + u_z\,\frac{\partial u_r}{\partial z}
    = \nu\,\frac{\partial^2 u_r}{\partial z^2},
    \label{eq:BLeqs}
\end{equation}
with boundary conditions
\begin{alignat}{2}
    u_r(r,0) &= 0 &&\qquad (\text{no-slip}),\label{eq:BC1}\\
    \frac{\partial u_r}{\partial z}\bigg|_{z=h(r)} &= 0 &&\qquad
    (\text{zero shear at free surface}),\label{eq:BC2}\\
    u_r(r,h) &= U_s = V_0 &&\qquad (\text{matching outer inviscid flow}),
    \label{eq:BC3}
\end{alignat}
where $h(r) = Q/(2\pi r V_0) = a^2 V_0/(2r V_0) = a^2/(2r)$
is the local film thickness (volume conservation, $a \equiv d/2$).

Watson introduces the similarity variable $\eta = z/\delta(r)$ and assumes
a cubic velocity profile within the boundary layer:
\begin{equation}
    \frac{u_r}{V_0} = \frac{3}{2}\eta - \frac{1}{2}\eta^3,
    \qquad 0 \leq \eta \leq 1,
    \label{eq:cubicprofile}
\end{equation}
which satisfies both~\eqref{eq:BC1} and~\eqref{eq:BC2}. The von K\'{a}rm\'{a}n
momentum integral, with outer surface velocity $U_s = V_0$ held constant
in Region~I (consistent with the inviscid outer film), reads
\begin{equation}
    \frac{39\,V_0^2}{280}\,\frac{d(r\,\delta)}{dr}
    = \frac{3\,r\,\nu\,V_0}{2\,\delta},
    \label{eq:karman}
\end{equation}
which admits the similarity solution $\delta^2 = (280/39)\,\nu\,r/V_0$, i.e.
\begin{equation}
    \frac{\delta(r)}{a} = c_\delta\,\sqrt{\frac{r}{a\,\Rey_a}},
    \qquad c_\delta = \sqrt{\tfrac{280}{39}} \approx 2.68,
    \label{eq:deltawatson}
\end{equation}
where $a = d/2$, $\Rey_a = V_0 a/\nu$. \emph{The Region~I boundary layer
grows as $\delta \propto r^{1/2}$, the standard Blasius scaling}, not as
$r^3$ (an algebraic mis-step in some informal derivations that yields a
dimensionally broken result).

The \emph{transition radius} $r_0$, at which the boundary layer spans the
entire film ($\delta = h$), follows from setting $\delta(r_0) = a^2/(2 r_0)$
in equation~\eqref{eq:deltawatson}:
\begin{equation}
    \frac{r_0}{a} = \left(\frac{\Rey_a}{4\,c_\delta^2}\right)^{1/3}
    \approx 0.330\,\Rey_a^{1/3},
    \label{eq:r0}
\end{equation}
in agreement (to within the cubic-profile error) with Watson's original
asymptotic $r_0/a \approx 0.316\,\Rey_a^{1/3}$~\citep{Watson1964}. For
$\Rey_a = V_0(d/2)/\nu = 7{,}300$ at $d = 1.3~mm$, $V_0 = 11.2~m/s$:
\begin{equation}
    r_0 = a\,(0.330)\,(7300)^{1/3} \approx a \times 6.4
    \approx 4.1~mm.
\end{equation}
This is much smaller than $r_\mathrm{max} \approx 15~mm$, so
\emph{the cleaning radius lies in Watson's Region~II, not Region~I}.
The Region~II shear law (derived in the next subsection) therefore
governs the cleaning constraint.

\subsubsection{Wall Shear Stress Distribution}

In Region~I ($r \leq r_0$), the wall shear from the cubic
profile~\eqref{eq:cubicprofile} is
\begin{equation}
    \tau_w(r) = \mu\left.\frac{\partial u_r}{\partial z}\right|_{z=0}
    = \frac{3\,\mu\,V_0}{2\,\delta(r)},
    \label{eq:taudef}
\end{equation}
and substituting~\eqref{eq:deltawatson} gives the corrected Region~I result
\begin{equation}
    \boxed{
    \tau_w(r) = \frac{3\,\rho V_0^2}{2\,c_\delta\,\Rey_r^{1/2}}
    \approx 0.560\,\rho V_0^2\,\Rey_r^{-1/2},
    \qquad r \leq r_0,
    }
    \label{eq:tauI}
\end{equation}
with $\Rey_r = V_0 r/\nu$. The Region~I wall shear scales as $r^{-1/2}$,
\emph{not} as $r^{-3}$. At $r = a$ this gives
$\tau_w(a) \approx 0.560\,\rho V_0^2/\sqrt{\Rey_a} \approx 650~Pa$
for $V_0 = 11.2~m/s$, $a = 0.65~mm$ --- a value of order the
jet dynamic pressure scaled by $\Rey_a^{-1/2}$, and physically
reasonable. (An earlier, incorrect derivation gave
$\tau_w(a) \approx 2.6\,\rho V_0^2 \approx 260~kPa$, an unphysical
overshoot by three orders of magnitude relative to direct measurements
from Phares et al.~\citep{Phares2000}, which find peak values of order
2--$10~kPa$ for jets in the relevant Reynolds-number range.)

Beyond $r_0$, the boundary layer spans the full viscous film. The
depth-averaged momentum balance with no streamwise pressure gradient,
parabolic film profile, and mass conservation $rh\bar u = Q/(2\pi)$
yields the asymptotic similarity solution $h(r) = (5/3)\,\nu\,r^2/(a^2 V_0)$
for $r \gg r_0$. The corresponding Region~II wall shear is
\begin{equation}
    \boxed{
    \tau_w(r) = c_2\,\rho V_0^2\,\Rey_a\left(\frac{a}{r}\right)^5,
    \qquad c_2 = 27/50 \approx 0.54,
    }
    \label{eq:tauII}
\end{equation}
valid for $r \gg r_0$. The decay is now $r^{-5}$, much steeper than
Region~I. \emph{For bidet parameters the cleaning radius lies in
Region~II}, so equation~\eqref{eq:tauII}, not~\eqref{eq:tauI}, is the
operative shear law for the cleaning criterion of
Section~\ref{sec:mass}.

Several features of the dual-region picture merit comment:
\begin{enumerate}[leftmargin=*, label=(\roman*)]
    \item \emph{No $r^{-3}$ pseudo-divergence at the origin.} The
    physical Region~I shear remains finite at small $r$
    (proportional to $\rho V_0^2/\sqrt{\Rey_r}$).
    \item \emph{Re-dependence is non-trivial.} In Region~II,
    $\tau_w \propto V_0^3 a^6/(\nu r^5)$: linear in fluid viscosity,
    cubic in jet velocity.
    \item \emph{Cleaning-zone average.} For $r$ in $[2a, r_\mathrm{max}]$
    spanning Region~I and Region~II, area-averaging
    \eqref{eq:tauI}/\eqref{eq:tauII} gives
    $\langle\tau_w\rangle \approx 50$--$500~Pa$, two orders of
    magnitude above $\tau_\mathrm{crit}$.
    \item \emph{Comparison to measurement.} Phares et al.~\citep{Phares2000}
    measure peak wall shear of $2$--$10~kPa$ for jets at $\Rey_d
    \sim 10^4$ and short standoff. The present analytical value at
    $r = a$ ($\sim$650~Pa) is below this peak because Watson's
    Region~I assumes laminar wall flow; turbulent transition would
    raise the peak by a factor of $3$--$10$, consistent with measurement.
    The cleaning-zone average is the relevant quantity, and the
    one to two orders of magnitude over $\tau_\mathrm{crit}$ is robust.
\end{enumerate}

\subsubsection{Hydraulic Jump}

At a radial distance $r_h$ determined by a balance between inertia,
viscosity, and gravity, the spreading film undergoes a \emph{circular
hydraulic jump}: an abrupt thickening accompanied by a rapid velocity
decrease. The horizontal-plate problem has a long history; the modern
shallow-water shock treatment of Bohr, Dimon \& Putkaradze
\citep{BohrEtAl1993} gives the asymptotic scaling
\begin{equation}
    r_h \;\sim\; \left(\frac{Q^5}{\nu^3\,g}\right)^{1/8}
    \;\propto\; Q^{5/8}\,\nu^{-3/8}\,g^{-1/8},
    \label{eq:hjump}
\end{equation}
later refined by surface-tension and downstream-depth corrections
\citep{BushAristoff2003}. For nominal bidet parameters
($Q \approx e-5~m^3/s$, $\nu = e-6~m^2/s$,
$g = 9.81~m/s^2$), \eqref{eq:hjump} gives $r_h$ of order tens of
millimetres, comparable to the anatomical boundary $r_\mathrm{max}$.
(The pure inviscid critical-Froude estimate $r_h \approx V_0^2 a^2/(2g)$
gives $r_h \sim 3~\mu  m$ for the same parameters and is not the
right scaling for a thin-film impinging jet.)
On the actual perianal anatomy the surface is concave and
oriented at a strongly non-horizontal angle relative to gravity, so
the flat-horizontal-plate jump is at best a guide; the principal
qualitative role of the jump is to mark the radius beyond which the
spreading film transitions to a slow, deep flow, providing a natural
soft boundary for the rinsing zone.

\subsection{Curvature Effects: Concave-Surface Perturbation}
\label{subsec:curvature}

We now generalize the Watson solution to the concave surface geometry.
Define the arc-length coordinate $s$ measured along the curved surface from
the stagnation point (replacing the flat-plate $r$). For $r \ll R$, we have
$s \approx r + r^3/(6R^2) + \ldots$, so to leading order $s \approx r$ and
the curvature correction enters at $\mathcal{O}(\epsilon^2)$.

The key dynamical modification on a concave surface is the centrifugal
pressure gradient in the direction normal to the surface ($n$-direction):
\begin{equation}
    \frac{\partial P}{\partial n} = -\frac{\rho U_e^2}{R_\mathrm{eff}},
    \label{eq:centrifugal}
\end{equation}
where $U_e$ is the outer edge velocity and $R_\mathrm{eff}(s)$ is the
local surface radius of curvature. Concave-wall flows are well known to be
\emph{destabilizing} in the sense of supporting the centrifugal
Taylor--G\"ortler instability (treated in the next subsection): the
centrifugal pressure gradient drives streamwise vortices that enhance
mixing and can amplify wall shear above the flat-plate baseline once
the instability is supercritical. Below the G\"ortler threshold, the
direct curvature correction to the laminar wall shear in equation~%
\eqref{eq:tauI}/\eqref{eq:tauII} is small:
\begin{equation}
    \tau_w^\mathrm{concave}(r) = \tau_w^\mathrm{flat}(r)
    \left[1 + \alpha_c\,\frac{\delta(r)}{R} + \mathcal{O}(\epsilon^2)\right],
    \label{eq:curvaturecorrection}
\end{equation}
with $\alpha_c = \order{1}$. For $\delta \approx 0.1~mm$ and
$R = 20~mm$, the first-order correction is $\delta/R \approx 0.005$,
a $\sim$0.5\,\% effect. Curvature is therefore a small correction in the
laminar regime, and a potentially \emph{shear-enhancing} correction once
G\"ortler vortices set in --- so the flat-plate Watson solution is a
conservative (under-)estimate of cleaning effectiveness on the concave
anatomy.

\subsubsection{Taylor--G\"ortler Vortex Instability}

Concave boundary layers are subject to the centrifugal \emph{Taylor--G\"ortler}
instability, which generates counter-rotating streamwise vortex pairs and
substantially enhances turbulent mixing and surface transport. The criterion
for onset is \citep{Hall1982}:
\begin{equation}
    \Go \equiv \frac{U_e\,\theta}{\nu}\,\sqrt{\frac{\theta}{R}} > \Go_\mathrm{crit}
    \approx 0.3,
    \label{eq:goertler}
\end{equation}
where $\theta$ is the momentum thickness of the boundary layer. For the
spreading water film at $r \approx r_0 \approx 9~mm$, the momentum
thickness $\theta \approx \delta/6 \approx 15~\mu  m$ and
$U_e \approx V_0$, giving $\Go \approx 0.03$, below the critical value,
meaning Taylor--G\"ortler vortices are unlikely in the thin wall-jet film
under typical bidet conditions. However, in Region~II ($r > r_0$), as the
film thickens and velocity decreases, $\Go$ grows and may approach critical
values near the hydraulic jump, where additional mixing would enhance late-stage
debris transport. This effect is conservative to neglect and can only improve
cleaning performance.

% =============================================================
\section{Mass Transfer and Cleaning Efficacy}
\label{sec:mass}
% =============================================================

\subsection{The Detachment Criterion}
\label{subsec:detach}

Cleaning of an adhered layer by fluid shear proceeds in two steps:
(i) detachment, when $\tau_w$ exceeds the adhesive yield stress; and
(ii) transport, when the dislodged material is carried away by the flow.
We address each in turn.

The critical wall shear stress $\tau_\mathrm{crit}$ for detachment of
fecal residue from skin is not directly measured in the literature (a gap
discussed in Section~\ref{sec:discussion}). We bound it from below and
above:

\begin{enumerate}[leftmargin=*, label=(\roman*)]
    \item \emph{Lower bound}: The EHEDG standard for industrial surface
    hygiene establishes $\tau_\mathrm{crit} = 3~Pa$ as the threshold
    for removing biological soil \citep{Ehedg2003}. Biofilm detachment
    studies find $\tau_\mathrm{crit} \approx 0.3$--$4.7~Pa$
    \citep{Teodosio2010}.
    \item \emph{Upper bound}: Fecal material has a measured dynamic yield
    stress of approximately $\sigma_y \approx 10~Pa$ at physiologically
    relevant moisture contents \citep{Woolley2014}. The skin-residue
    adhesion is plausibly weaker still, since the contact involves a
    moist mucosal interface rather than a fully consolidated bolus, but
    we are unaware of a direct measurement on perianal substrates ---
    a gap discussed in Section~\ref{sec:discussion}.
\end{enumerate}

We therefore adopt $\tau_\mathrm{crit} = 1~$--$10~Pa$ as
the operative range. Since the cleaning radius lies in Region~II for
typical bidet parameters (Section~\ref{subsec:watson}), we set
$\tau_w(r_\mathrm{clean}) = \tau_\mathrm{crit}$ in
equation~\eqref{eq:tauII}:
\begin{equation}
    r_\mathrm{clean} = a\left(\frac{c_2\,\rho V_0^2\,\Rey_a}{\tau_\mathrm{crit}}\right)^{1/5}
    = \left(\frac{c_2\,\rho V_0^3 a^6}{\tau_\mathrm{crit}\,\nu}\right)^{1/5}.
    \label{eq:rclean}
\end{equation}
For $\tau_\mathrm{crit} = 5~Pa$, $V_0 = 11.2~m/s$,
$a = 0.65~mm$, $\nu = e-6~m^2/s$:
\begin{align}
    r_\mathrm{clean} &= \left(\frac{0.54 \times 1000 \times (11.2)^3
    \times (6.5\times 10^{-4})^6}{5\,\mathrm{Pa}\times 10^{-6}\,\mathrm{m^2/s}}\right)^{1/5}
    \nonumber\\
    &\approx 26~mm\quad
    (\text{Region~II laminar lower bound, $H/d \to 1$}).
    \label{eq:rcleanval}
\end{align}
(The earlier draft reported $\approx 16$~mm; the corrected arithmetic
is given here.) For the realistic seated-user standoff $H/d \approx 60$,
the developed-jet decay $V_c/V_0 = B d/H \approx 0.1$ reduces the
effective driving velocity by $\sim 10\times$ and, via the
$V_0^{3/5}$ scaling of \eqref{eq:rclean}, the laminar
$r_\mathrm{clean}$ to $\approx 26 \times 0.1^{3/5} \approx 6.5~mm$.
Turbulent wall-jet enhancement at $\Rey_d \sim 10^4$ raises this back
above $r_\mathrm{max}$ in practice (see Sec.~\ref{subsec:peclet}
discussion and the companion optimization
analysis~\citep{BurtonOptPaper}), and a sweep through $\pm 15^\circ$
covers the unreached annulus directly.

The corrected $r_\mathrm{clean}$ ($\approx 26~mm$ in the
close-deploy laminar limit) more than encloses the anatomical
contamination zone ($r_\mathrm{max} \approx 10$--$20~mm$); the
realistic seated-user value sits at the lower edge of $r_\mathrm{max}$,
and an oscillation lever or close-deploy architecture provides
robust safety margin. Since $r_\mathrm{clean} \gg r_0$, cleaning
occurs primarily in Watson's Region~II, where the $r^{-5}$ shear stress
decays rapidly with $r$; this means the cleaning radius is geometrically
bounded and a small change in $V_0$ or $a$ produces a strong response in
$r_\mathrm{clean}$ ($\propto V_0^{3/5} a^{6/5}$).

The margin of safety is defined as
\begin{equation}
    \mathcal{M} = \frac{\langle\tau_w\rangle}{\tau_\mathrm{crit}} \approx
    \frac{50\text{--}500\,\mathrm{Pa}}{1\text{--}10\,\mathrm{Pa}} = 5\text{--}500,
    \label{eq:margin}
\end{equation}
a safety margin of one to nearly three orders of magnitude.

\subsection{Advective Removal: The P\'{e}clet Number}
\label{subsec:peclet}

Detachment is necessary but not sufficient for cleaning: dislodged material
must be transported away from the surface. The competition between advective
and diffusive transport is quantified by the P\'{e}clet number:
\begin{equation}
    \Pe = \frac{U_\mathrm{film}\,L}{\mathcal{D}},
    \label{eq:peclet}
\end{equation}
where $U_\mathrm{film} \approx 2~m/s$ is a representative surface
velocity in the cleaning zone (the Region~II surface velocity
$U_s(r) = (3/2)\bar u \propto V_0^2 a^4/(\nu r^3)$ evaluated at
$r \approx r_\mathrm{clean}/2$, with $V_0 \to V_c$ for typical standoff),
$L \approx r_\mathrm{clean} \approx 0.016~m$,
and $\mathcal{D} \sim 10^{-10}$--$10^{-9}\,\mathrm{m^2/s}$ is the effective
diffusivity of fecal particulates (micron-scale; Stokes--Einstein).
\begin{equation}
    \Pe \sim \frac{(2)(0.016)}{10^{-10}} = 3.2 \times 10^8,
    \label{eq:peval}
\end{equation}
i.e.\ $\Pe \sim 10^{8}$--$10^9$. Even allowing for two-decade
uncertainty, transport is overwhelmingly advective. Once detached, debris particles are
swept radially outward by the spreading film in a time $t \sim L/U_\mathrm{film}
\approx 3.4~ms$, before falling under gravity into the toilet bowl.
This single-pass removal mechanism is the central advantage of water over
paper: toilet paper displaces fecal material but provides no transport fluid,
leading to smearing rather than removal. A quantitative estimate of residual
contamination after wiping is beyond the scope of this paper (and beyond
the scope of any published clinical study; see Section~\ref{sec:discussion});
however, the absence of any transport mechanism in dry wiping is physically
unambiguous.

% =============================================================
\section{Tissue Mechanics: The Cost of Wiping}
\label{sec:tribology}
% =============================================================

\subsection{Tribological Analysis}
\label{subsec:tribology}

The mechanical interaction between toilet paper and perianal skin is governed
by tribological principles. The friction force during a wiping pass is
\begin{equation}
    F_f = \mu_k\,N,
    \label{eq:friction}
\end{equation}
where $N$ is the normal load and $\mu_k$ is the kinetic friction coefficient.
The resulting contact shear stress is
\begin{equation}
    \tau_\mathrm{wipe} = \frac{F_f}{A_c} = \mu_k\,P_c,
    \label{eq:contactshear}
\end{equation}
where $P_c = N/A_c$ is the contact pressure.

Derler and Gerhardt \citep{Derler2012} compiled extensive data on skin
friction, reporting $\mu_k = 0.2$--$2.7$ with strong dependence on
moisture, temperature, and substrate material. On moist skin (relevant to
post-defecation conditions), paper-on-skin friction coefficients of
$\mu_k = 0.5$--$1.5$ are well documented \citep{Vilhena2023}; wet toilet
paper on wet skin can reach $\mu_k \approx 2.0$--$4.1$ due to meniscus
adhesion. We adopt the conservative range $\mu_k = 0.5$--$1.5$.

Applied wiping forces for perianal hygiene in normal adults have not been
directly published. From ergonomic and pressure-sensor data for analogous
manual tasks (clinical perineal care), we estimate $N \approx 2$--$10~N$
applied over a contact area $A_c \approx 5$--$15~cm^2$ (a folded pad
of 4--6 ply), giving contact pressures $P_c \approx 2$--$20~kPa$.
The wiping shear stress is therefore:
\begin{equation}
    \tau_\mathrm{wipe} = \mu_k\,P_c \approx (0.5\text{--}1.5)
    \times (2{,}000\text{--}20{,}000\,\mathrm{Pa})
    = 1{,}000\text{--}30{,}000\,\mathrm{Pa}.
    \label{eq:wiperange}
\end{equation}

\subsection{Quantitative Stress Comparison}

Table~\ref{tab:stresscomparison} compares the mechanical stresses delivered
by each modality. The bidet's hydrodynamic wall shear occupies the range
1--100~Pa (averages over the cleaning zone), while wiping
generates 1--30~kPa, a ratio of two to three orders of
magnitude.

\begin{table}[h!]
\centering
\caption{Comparative mechanical stress analysis: bidet vs.\ toilet paper.}
\label{tab:stresscomparison}
\begin{tabular}{lccc}
\toprule
\textbf{Parameter} & \textbf{Bidet (water jet)} & \textbf{Toilet paper} & \textbf{Ratio (TP/Bidet)}\\
\midrule
Shear stress on tissue     & $1$--$100~Pa$    & 1000~Pa--30000~Pa & $10$--$1000\times$\\
Normal contact pressure    & $0.4$--$2~kPa$ (stagnation) & $2$--$20~kPa$ & $5$--$50\times$\\
Contact type               & Non-contact (hydrodynamic) & Solid friction (abrasive) & N/A\\
Transport mechanism        & Advective ($\Pe \sim 10^9$)   & None (smearing)         & N/A\\
Friction coefficient       & $0$ (no solid contact) & $\mu_k = 0.5$--$1.5$   & N/A\\
\bottomrule
\end{tabular}
\end{table}

\subsection{Stratum Corneum Disruption}
\label{subsec:SC}

The perianal stratum corneum (SC) is the outermost, rate-limiting barrier
of the skin, comprising 10--20 corneocyte layers at a total thickness of
$\sim15~\mu  m$ \citep{Loden2012}. The SC is more susceptible to
friction injury than skin at other body sites due to the combined effects
of chronic moisture exposure (occlusion reduces barrier function) and thin
tissue architecture.

Each wiping pass applies the shear stress in equation~\eqref{eq:wiperange}
over the SC, progressively disrupting the intercorneocyte lipid lamellae.
L\"offler \textit{et al.}\ \citep{Loffler2004} demonstrated that barrier
disruption (measured by trans-epidermal water loss, TEWL) is directly
proportional to both the frictional force and the number of passes.
For 3--6 passes per event and multiple events per day, cumulative SC
disruption is substantial.

The clinical injury threshold for \emph{pressure injury} formation (a more
severe endpoint than barrier disruption) is approximately
$\tau_\mathrm{injury} \approx 9$--$15~kPa$ \citep{Broderick2021}.
Vigorous wiping ($\tau_\mathrm{wipe} \approx 10~kPa$) reaches this
threshold, particularly in elderly or immunocompromised patients.

The bidet delivers no solid-contact stress. The hydrodynamic wall shear
($\tau_w \approx 10$--$100~Pa$) is $100$--$1{,}000\times$ below the
SC disruption threshold, and no friction coefficient enters the problem.
The difference in tissue loading between the two modalities is not
incremental but physical in kind.

\subsection{Clinical Sequelae}
\label{subsec:clinical}

The mechanical analysis predicts the following hierarchy of tissue pathology
from chronic toilet-paper use, consistent with the clinical literature:

\begin{enumerate}[leftmargin=*]
    \item \emph{Pruritus ani}: Affects 1--5\% of the general population.
    Dasan \textit{et al.}\ \citep{Dasan1999} found contact dermatitis (including
    from toilet paper allergens: formaldehyde, chlorine residues, fragrances)
    in 53\% of identifiable cases at a specialized clinic.
    \item \emph{Hemorrhoidal exacerbation}: Hemorrhoids affect $\sim10$ million
    Americans annually \citep{Peery2014}. Ryoo \textit{et al.}\ \citep{Ryoo2011}
    demonstrated that bidet use at low jet forces (40--200~mN) significantly
    reduces anal resting pressure, whereas wiping with paper at 2--10~N
    exacerbates vascular engorgement.
    \item \emph{Anal fissures}: Chronic SC disruption reduces barrier function
    against fecal bacteria and increases susceptibility to fissure formation,
    particularly when combined with constipation-induced trauma.
\end{enumerate}

The systematic review by Baig \textit{et al.}\ \citep{Baig2022} should be
interpreted with caution: their finding of elevated bidet use among subjects
with perianal pathology most likely reflects reverse causation (patients adopt
bidets \emph{because} of symptoms), confirmed by the prospective study of
Kiuchi \textit{et al.}\ \citep{Kiuchi2017}, which found no significant incidence
odds ratio for new hemorrhoid development among bidet users.

% =============================================================
\section{Environmental and Economic Analysis}
\label{sec:enviro}
% =============================================================

\subsection{Resource Consumption of the American Toilet Paper Habit}

At an estimated $\sim$141~rolls/ person/ year consumption rate
applied to a $\sim 330$ million population, the United States consumes
approximately $46$ billion rolls of toilet paper annually, representing
the highest per-capita usage on Earth and a disproportionate share of
the global supply~\citep{ScientificAmerican2009}. The environmental
footprint per roll is estimated at: 37 gallons of embedded water (tree
growth and manufacturing), 0.47 kWh of electricity, and 140 g of
CO$_2$-equivalent emissions (virgin-fiber product) \citep{NRDC2020}.
Aggregated to national scale:
\begin{itemize}[leftmargin=*]
    \item Embedded water: $\sim$1.7~trillion gal/yr
    ($46 \times 10^9 \times 37~gal \approx
    1.72 \times 10^{12}\,\mathrm{gal/yr}$);
    \item Electricity: $\sim$22~TWh/yr
    ($46 \times 10^9 \times 0.47~kWh$);
    \item Virgin trees harvested: $\sim$15--31 million annually, sourced
    disproportionately from Canada's boreal forest;
    \item Carbon emissions from production:
    $\sim$6.5~Mt\,CO_2\text{-eq/yr}
    ($46 \times 10^9 \times 140\,\mathrm{g} \approx
    6.4 \times 10^9\,\mathrm{kg/yr}$); lifecycle figures including
    upstream forest disturbance and downstream wastewater treatment
    are typically 2--3$\times$ higher.
\end{itemize}

\subsection{Lifecycle Comparison}

A single bidet use consumes $Q_\mathrm{bidet} \approx 0.5~L$ of water.
For 2 uses per day per person:
\begin{equation}
    W_\mathrm{bidet} = 0.5\,\mathrm{L} \times 2 \times 365 = 365\,\mathrm{L/yr}
    \approx 96~gal/yr.
\end{equation}
Compared to the $5{,}200$\,gal/yr in embedded TP water, the bidet uses
$\sim\!1.8\%$ of the water budget. The net embedded-water reduction from
national adoption would be $\sim\!5{,}100\,\mathrm{gal\,yr^{-1}\,person^{-1}}$,
or $\sim\!1.7$ trillion gallons nationally.

Non-electric bidet attachments carry a manufacturing carbon footprint of
$\sim\!3$--$6\,\mathrm{lbs\,CO_2e}$ (embodied in $\sim\!200\,\mathrm{g}$
of plastic), amortized over a 5--10 year lifespan: negligible compared to
the annual avoided carbon footprint of $\sim\!50\,\mathrm{lbs\,CO_2e\,yr^{-1}}$
from eliminated TP production.

\subsection{Economic Analysis}

Define household variables: $n$ occupants, $C_\mathrm{TP}$ = annual toilet
paper expenditure, $C_\mathrm{bidet}$ = bidet purchase price, $f$ = fractional
TP reduction post-adoption (typically $f = 0.65$--$0.75$), $r$ = real discount
rate. The \emph{net present value} of bidet adoption over a horizon $T$ is:
\begin{equation}
    \mathrm{NPV} = -C_\mathrm{bidet} + \sum_{t=1}^{T}\frac{f\,C_\mathrm{TP}}{(1+r)^t}
    = -C_\mathrm{bidet} + f\,C_\mathrm{TP}\,\frac{1-(1+r)^{-T}}{r}.
    \label{eq:npv}
\end{equation}

For $C_\mathrm{bidet} = \$50$ (attachment), $C_\mathrm{TP} = \$200/\mathrm{yr}$
(4-person household), $f = 0.70$, $r = 0.05$, $T = 10\,\mathrm{yr}$:
\begin{equation}
    \mathrm{NPV} = -50 + (0.70)(200)\,(7.72) = -50 + 1{,}081 = \$1{,}031.
    \label{eq:npvval}
\end{equation}

The \emph{payback period} $T^*$ is defined by $\mathrm{NPV}(T^*) = 0$:
\begin{equation}
    T^* = \frac{C_\mathrm{bidet}}{f\,C_\mathrm{TP}}
    = \frac{50}{(0.70)(200)} \approx 0.36\,\mathrm{yr} \approx
    4.3~months.
    \label{eq:payback}
\end{equation}

At the national scale, if 50\% of 130 million US households adopted a
bidet attachment:
\begin{equation}
    \mathrm{Annual\;savings} = 0.50 \times 1.3 \times 10^8 \times \$140
    \approx \$9.1\,\mathrm{billion/yr}.
\end{equation}

Additional healthcare cost offsets, from reduced incidence of toilet-paper-exacerbated
perianal pathology, represent a speculative but plausibly large additional
benefit: US hemorrhoid-related healthcare costs alone exceed \$770 million
per year (employer-insured population only) \citep{Peery2014}, with the OTC
hemorrhoid treatment market adding a further \$850 million in North America.

% =============================================================
\section{Policy Framework and Recommendations}
\label{sec:policy}
% =============================================================

The United States currently has no federal or state incentive program for
bidet adoption. The EPA WaterSense program certifies low-flow toilets and
provides a well-established administrative framework for bathroom fixture
rebates (typical value: \$50--\$200 from participating utilities). We propose
the following policy interventions, in increasing order of regulatory ambition:

\begin{enumerate}[leftmargin=*, label=\textbf{P\arabic*.}]
    \item \textbf{WaterSense Bidet Certification}: Expand the WaterSense
    label to encompass bidet attachments and seats meeting a flow threshold
    of $\leq0.5~L/ use$. This requires no new legislation and
    could be implemented through administrative rulemaking.
    \item \textbf{Utility Rebate Programs}: Extend existing WaterSense
    rebate frameworks to include qualifying bidet attachments at \$25--\$50
    per unit, funded through water utility conservation budgets. The payback
    period on utility investment is immediate given the embedded-water savings.
    \item \textbf{FSA/HSA Blanket Eligibility}: Extend Flexible Spending
    Account and Health Savings Account eligibility for bidet attachments
    to all users (currently limited to those with documented medical need),
    reducing effective consumer price by 20--37\% through tax pre-emption.
    \item \textbf{Building Code Update}: Require GFCI-protected electrical
    outlets within 1~m of toilets in all new residential construction.
    This low-cost building code amendment (estimated \$25--\$75 per unit in
    new construction) removes the primary electrical installation barrier
    for electric bidet seats without mandating bidet installation itself.
    \item \textbf{Public Procurement Mandate}: Require bidet installation
    in all new federal buildings, VA facilities, and federally funded
    healthcare infrastructure, a market-size intervention analogous to
    low-flow toilet mandates of the 1990s.
\end{enumerate}

Japan's trajectory (14\% adoption in 1992 to 80\% in 2018) demonstrates
that the transition is achievable on a 25-year timescale without regulatory
mandates, driven primarily by consumer education, cost reduction, and
infrastructure normalization \citep{JapaneseCabinet2018}. The US starting
position ($\sim3\%$ adoption) is analogous to Japan circa 1985.

% =============================================================
\section{Discussion: Evidentiary Gaps}
\label{sec:discussion}
% =============================================================

The following gaps remain in the current evidence base:

\begin{enumerate}[leftmargin=*, label=(\roman*)]
    \item \emph{Fecal adhesion strength to skin}: No published study directly
    measures the critical shear stress for detachment of fecal material from
    perianal skin. Our estimate of $\tau_\mathrm{crit} = 1$--10~Pa is
    extrapolated from industrial biofilm literature and fecal rheology data.
    A dedicated in-vitro experiment using representative fecal simulants on
    excised or synthetic skin would substantially tighten this range.
    \item \emph{RCT on cleaning efficacy}: No large-scale randomized
    controlled trial has directly compared residual fecal contamination
    (colony counts or chemical biomarkers) between bidet and toilet-paper
    cohorts under controlled conditions. This is the single most important
    clinical evidence gap. The 2005 \textit{Gastroenterology Nursing} study
    is the closest analog but lacked adequate sample size and blinding.
    \item \emph{Wiping force measurement}: No published study reports direct
    force-plate measurements of applied normal load during toilet-paper wiping.
    Our range of 2--10~N is inferred from ergonomic analogs. Instrumented
    toilet paper (force-sensitive substrate) would be straightforward to
    implement.
    \item \emph{National bidet adoption statistics}: No US government survey
    tracks bidet household penetration. The $\sim3\%$ figure is an industry
    estimate; a Census or NHANES supplement question would enable population-level
    epidemiological analyses linking adoption to perianal health outcomes.
    \item \emph{Heterogeneity of bidet jet design}: This analysis uses
    nominal parameters for a standard non-oscillating jet. Manufacturers
    such as TOTO employ oscillating jets, varied spray angles (43$^\circ$ rear,
    53$^\circ$ front), and adjustable flow rates. A parametric analysis of
    cleaning efficacy across the manufacturer design space is warranted.
    \item \emph{Laminar Region~II conservatism}: The wall-shear
    estimates here use Watson's laminar similarity asymptotic, which
    is a lower bound on $\tau_w$ at the operative
    $\Rey_d \sim 10^4$. The Phares et al.~\citep{Phares2000}
    correlation implies that the actual turbulent wall jet decays as
    $\tau_w \propto r^{-1.05}$ rather than $r^{-5}$, giving wall
    shear several-fold higher than the present estimates. This
    tightens, not weakens, the case for hydrodynamic cleansing.
\end{enumerate}

The companion optimization manuscript~\citep{BurtonOptPaper} scopes a
workstation-scale verification programme --- reduced-order numerical
Watson solver, axisymmetric VOF CFD on flat and concave plates,
oblique/oscillating-jet RANS, and a benchtop PIV / hot-film rig with a
microfluidic shear cell for direct $\tau_\mathrm{crit}$ measurement ---
in four effort tiers suitable for a Python laptop through to a
$\sim$$\$$$1500$ rig build. The most efficient path to closing the gaps
listed above is a $\sim$one-week CFD subset (flat plate, single
$H/d, \Rey_d$) that converts the laminar-Region~II conservatism into a
quantitative turbulent enhancement factor, followed by the microfluidic
$\tau_\mathrm{crit}$ measurement.

% =============================================================
\section{Conclusions}
\label{sec:conclusion}
% =============================================================

The core quantitative results are:

\begin{enumerate}[leftmargin=*]
    \item The bidet jet ($\Rey \approx 14{,}000$, $\We \approx 2{,}000$)
    arrives at the skin surface as a coherent liquid column, generating
    wall shear stresses of $\sim$50--500\,Pa across the cleaning zone, one
    to nearly three orders of magnitude above the critical detachment
    threshold of 1--10\,Pa for fecal material.
    \item The corrected Watson two-region wall-jet analysis predicts an
    effective cleaning radius $r_\mathrm{clean} \approx 16~mm$,
    comfortably enclosing the contaminated perianal zone, with a safety
    margin $\mathcal{M} \approx 5$--$500$.
    \item Concave surface curvature ($\epsilon = d/R \approx 0.05$--$0.10$)
    introduces a small first-order shear correction; in the supercritical
    G\"ortler regime, streamwise vortices would further enhance late-stage
    debris transport near the hydraulic jump.
    \item The P\'{e}clet number $\Pe \sim 10^{8}$--$10^{9}$ confirms that
    transport is overwhelmingly advective: dislodged debris is removed in
    $\sim$8~ms by the spreading film.
    \item Toilet paper wiping delivers shear stresses of 1{,}000--30{,}000\,Pa
    to moist perianal skin, two to three orders of magnitude above the
    bidet's hydrodynamic load, with no transport mechanism, generating
    progressive stratum corneum disruption and exacerbating a clinical
    spectrum of perianal pathology.
    \item National bidet adoption would eliminate 15--31 million tree
    harvests annually, reduce embedded water consumption by
    $\sim$98~\% (saving $\sim$1.7~trillion gal/yr), and
    deliver net present values of $>\$1{,}000$ per household over 10 years
    at a payback period of $\sim$4 months.
\end{enumerate}

The primary barrier to adoption in the United States is infrastructure
(electrical outlets near toilets) and cultural inertia, not engineering,
cost, or efficacy. We propose an EPA WaterSense certification expansion,
utility rebate programs, and a building code amendment requiring GFCI outlets
at toilet locations as the minimum necessary policy package.

% =============================================================
%